
\section{Taškinės grupės}

Taškinės grupės yra grupės $\Orth3$ pogrupiai.
Jas galime skirstyti pagal tai kokius $\set R^3$ taškus jos fiksuoja.
Taip gauname septynias begalines ašinių grupių šeimas
ir septynias poliedrines grupes. 
Taškines grupes taip pat galime skirstyti pagal tai ar jos yra $\SO3$ pogrupiai.
Taškinės grupės, kurios yra $\SO3$ pogrupiai, yra vadinamos chiralinėmis grupėmis,
Taškinės grupės, kurioms priklauso netikrieji sukimai ir todėl jos nėra $\SO3$ pogrupiai,
yra vadinamos achiralinėmis grupėmis.
Visas šias grupes apibrėšime pateikdami jų generatorius ir sąryšius,
kuriuos šie generatoriai privalo tenkinti. 



\subsection{Grupių prezentacijos}

Grupės prazentacija yra būdas apibrėžti grupę,
pateikiant jos generatorius ir sąryšius kuriuos jie tenkina.
Formaliai tokia grupė yra faktor-grupė laisvosios grupės,
generuojamos duotų generatorių.
Toliau pristatysime kobinatorinės grupių teorijos pagrindus,
reikalingus apibrėžti grupės prezentacija.

Tegul $S = \qty{a,b,c,\cdots}$ yra aibė generatorių.
Literatūroje įprasta aibę $S$ vadinti \emph{abėcėle}, 
o jos elementus - \emph{raidėmis}.
Tegul $S^{-1} = \qty{a^{-1},b^{-1},c^{-1},\cdots}$
yra aibė sudaryta iš \emph{atvikštinių raidžių}
(kol kas šios raidės yra tik aibės $S^{-1}$ elementai be papildomos struktūros).
Aibę $S \cup S^{-1}$ vadinsime \emph{išplėstine abėcėle}.
Baigtines sekas $ f_1 f_2 \cdots f_n$ vadinsime \emph{žodžiais},
čia visi $f_i \in S \cup S^{-1}$.
Tegul $e$ žymi tuščia žodį,
tai yra žodį sudarytą iš nulio raidžių.
Tegul tarp dviejų žodžių yra ekvivalntiškumo sąryšis
$f_1 f_2 \cdots f_n \sim f'_1 f'_2 \cdots f'_m$,
jei iš vieno žodžio galime gauti kitą
pakartotinai įterpdami raidės ir jos anvikštinės raidė poras $xx^{-1}$ ir $x^{-1}x$, $x \in S$.
Visi žodžiai tarp kurių yra ekvivalentiškumo sąryšis sudaro ekivalentiškumo klasę.
\begin{definition}
	Laisvoji grupė $F_S$ generuojama \emph{abėcėlės} $S$ yra sudaryta iš
	\emph{žodžių} ekvivalentiškumo klasių,
	kartu su daugyba duodama \emph{žodžių} (klasės atstovų) sujungimo.
\end{definition}

\begin{example}
	Parodykime kad $F_S$ tenkina grupės aksiomas.
	\begin{itemize}
		\item Uždarumas:
			sujungus du žodžius gausime nauja baigtinę raidžių seką, kuri taip pat yra žodis.
		\item Asociatyvumas:
			sekų jungimas yra asociatyvus.
		\item Egzistuoja vienetinis elementas:
			betkurį žodžį sujungus su tuščiu žodžiu $e$, gausime nepakeistą žodį.
		\item Egzistuoja atvikštiniai elementai:
			žodžiui $f_1 f_2 \cdots f_n$ atvikštinis yra $f_{n}^{-1} f_{n-1}^{-1} \cdots f_1^{-1}$,
			kur 
			\begin{equation}
				f_i^{-1}= 
				\begin{cases}
					x^{-1}\,,\;	&\text{kai } f_i=x
					\\
					x\,,		&\text{kai } f_i=x^{-1}
				\end{cases}
				\,,\qquad
				x \in S
				\,.
			\end{equation}
			Šių žodžių sandauga priklauso tai pačiai ekvivalentiškumo klasei,
			kaip tuščias žodis $e$.
			\begin{equation}
				f_1 f_2 \cdots f_n f_{n}^{-1} f_{n-1}^{-1} \cdots f_1^{-1}
				\sim
				f_1 f_2 \cdots f_{n-1} f_{n-1}^{-1} f_{n-2}^{-1} \cdots f_1^{-1}
				\sim \cdots \sim 
				f_1 f_1^{-1} \sim e
			\end{equation}
	\end{itemize}
\end{example}

Laisvoji grupė $F_S$ yra bendriausia įmanoma grupė turinti elementus $\qty{a, b, c, \cdots}$.
Ją sudarėme visų žodžių aibei suteikdami sąryšio $xx^{-1}=e$ stuktūra ir asociatyvia uždara daugybą,
tai yra minimali struktūra reikalinga kad $F_S$ būtų grupė.
Toliau parodysime, kaip galima sudaryti šios grupės faktor-grupes,
kurios turi papildomą struktūrą.

Yra aibė $R \subset F_S$,
jos elementai vadinami \emph{saryšio žodžiais} (ang. \emph{relators}).
Tai bus žodžiai, kurie faktor-grupėje bus sutapatinti su vienetiniu elementu.
Tegul $N(R)$ yra aibės $R$ normalusis uždarinys (ang. \emph{normal closure})
\begin{equation}
	N(R) := \qty{
		\left.
		\prod_{i=1}^n g_i r_i g_i^{-1}
		\;\right|\;
		n \in \set N ,\,
		r_i \in R \cup R^{-1} ,\,
		g_i \in F_S
	}
	\,.
\end{equation}

\begin{proposition}
	$N(R)$ yra grupės $F_S$ normalusis pogrupis.
\end{proposition}
\begin{proof}
	Parodome pogrupio uždarumą.
	$N(R)$ yra visų sąryšio žodžių jungtinių elementų baigtinės sandaugos.
	Sudauginus dvi tokias sandaugas gaunama nauja baigtinė sandauga,
	tad $N(R)$ yra uždara daugybos atžvilgiu.

	Parodome, kad $gN(R)g^{-1} \subseteq N(R)$ su visais $g \in F_S$.
	\begin{equation}
		g\qty(\prod_{i=1}^n g_i r_i g_i^{-1})g^{-1} 
		= \prod_{i=1}^n g g_i r_i g_i^{-1}g^{-1}
		= \prod_{i=1}^n g g_i r_i (g g_i)^{-1}
		\,,
	\end{equation}
	elementai $g g_i$ priklauso grupei $F_S$,
	tad pagal $N(R)$ apibrėžimą $g\qty(\prod_{i=1}^n g_i r_i g_i^{-1})g^{-1} \in N(R)$.
\end{proof}


\begin{definition}
	Turint \emph{abėcėlę} $S$ ir \emph{sąryšio žodžius} $R$,
	grupės prezentacija $\present{S}{R}$ yra faktor-grupė $F_S/N(R)$.
\end{definition}

\begin{remark}
	Literatūroje yra dažniau naudojamas žymėjimas $\langle\,S\,|\,R\,\rangle$,
	bet mes pasirinkome kitokį žymėjimą,
	kad nekiltų rizikos sumaišyti grupės prezentacija su Dirac'o ,,bra-ket'' žymėjimu.
	Taip pat abėcėlė $S$ bus pateikta nekaip aibė, o elementais.
	Sąryšio žodžiai $R$ bus pateikti kaip grupę apibrėžiantys sąryšiai
	$r' = r$, kur $r'r^{-1}\in R$.
	Sąryšiai yra patogūs, nes $F_s$ elemetai $r$ ir $r'$ priklauso
	tai pačiai gretutinei klasei
	\begin{equation}
		r\, N(R)
		= r\, r^{-1}	(r'r^{-1})	r \,N(R)
		= r'\,N(R)
		\,,
	\end{equation}
	tad grupėje $\present{S}{\qty{r'r^{-1},\cdots}}$
	turime tapatybę $r' = r$ (rašome tik gretutinės klasės atstovus).
\end{remark}

\begin{example}
	Grupės prezentaciją
	$ \present{S}{R} $, $ S = \qty{a, b} $,
	$ R = \qty{a^3, b^2, aba^{-1}b^{-1}} $,
	rašysime, taip 
	\begin{equation}
		\present{a,b}{a^3 = b^2 = e ,\, ab = ba}
	\end{equation}
	ir su grupės elelementais galime atlikti skaičiavimus kaip kad
	\begin{equation}
		abab = abba = aea = a^2
		\,.
	\end{equation}
\end{example}


\begin{remark}
	Svarbu paminėti, kad grupės presentacija nėra unikali.
	\begin{itemize}
		\item Galime pridėti sąryšio žodį $w \in N(R)$, $w \notin R$
			ir gausime lygias grupių prezentacijas
			\begin{equation}
				\present S R = \present{S}{R \cup \qty{w}}
				\,.
			\end{equation}
		\item Jei buvo pasirinkta daugiau sąryšio žodžių negu būtina,
			galima dalį sąryšio žodžių išmesti iš $R$.
		\item Galime pridėti generatorių $d$ prie abėcėlės $S$
			ir naują sąryšio žodį $dw^{-1}$ prie $R$,
			kur $w$ yra bet koks žodis iš piminės abėcėlės $S$ raidžių,
			ir gausime izomorfiškas grupes
			\begin{equation}
				\present S R \cong \present{ S \cup \qty{d} }{ R \cup \qty{dw^{-1}} }
				\,.
			\end{equation}
		\item Jei buvo pasirinkta daugiau generatorių negu būtina,
			galima dalį generatorių pakeisti žodžiais,
			sudarytais iš likusių generatorių.
	\end{itemize}
	Šie grupių prezentacijų pakeitimai vadinami Tieze transformacijomis.
	Klausimai, ar dvi duotos grupių prezentacijos yra izomorfiškos,
	ir ar duota grupės prezentacija yra minimali
	(neturi generatorių ar sąryšio žodžių,
	kuriuos išmetus gautume izomorfišką grupę), -
	yra netrivialūs kombinatorinės grupių teorijos uždaviniai.
\end{remark}

\begin{example}
	Parodysime, kad grupės prezentacija $\present{a,b}{a^2=b^2=ab=e}$
	yra izomorfiška trivialiai grupei $E$.
	\begin{equation}
		\begin{aligned}
			ab=e	&\;\;\text{ir}\;\; a^2=e	&\implies&& a &= b		\,,\\
			a=b		&\;\;\text{ir}\;\; b^3=e	&\implies&& a^3 &= e	\,,\\
			a^3=e	&\;\;\text{ir}\;\; a^2=e	&\implies&& a &= e		\,,\\
			a=e		&\;\;\text{ir}\;\; a=b		&\implies&& b &= e		\,.
		\end{aligned}
	\end{equation}
	Abu generatoriai yra lygūs vienetiniam elementui,
	tad $ \present{a,b}{a^2=b^2=ab=e} \cong E $.
\end{example}




\subsection{Ašinės grupės}

Ašinėmis grupėmis yra vadinami $\Orth3$ pogrupiai,
kurie veikdami į $\set R^3$ fiksuoja taškus priklausančius $z$ ašiai.
Ašinės grupės yra skirstomos į septynias begalines šeimas.

\begin{definition}
	Ciklinės grupės yra
	\begin{equation}
		C_n := \present{r}{r^n=e}
		\,,\quad
		n \in \set N
		\,,
	\end{equation}
	čia generatorius $r$ yra tapatinamas 
	su sukimu kampu $2\pi/n$ aplink $z$ ašį.
\end{definition}

\begin{definition}
	Diiedrinės grupės (ang. \emph{dihedral groups}) yra
	\begin{equation}
		D_n := \present{r,s}{r^n=s^2=e,\, srs^{-1}=r^{-1}}
		\,,\quad
		n \in \set N
		\,,
	\end{equation}
	čia generatorius $r$ - vėl tapatinamas su $2\pi/n$ pasukimu aplink ašį $z$,
	generatorius $s$ yra tapatinamas sukimu kampu $\pi$ aplink ašį statmeną ašiai $z$.
\end{definition}

\begin{proposition}
	Grupė $D_n$ yra izomorphiška pusiau tiesioginei grupių sandaugai $C_n \ltimes C_2$.
	Tei ekvivalentu teiginiui, kad egzistuoja trumpoti tikslioji seka
	\begin{equation}
		\begin{tikzcd}
			1 \arrow[r] &C_n \arrow[r,"\zeta"]	&D_n \arrow[r, "\xi"] &C_2 \arrow[r] &1
		\end{tikzcd}
	\end{equation}
	ir ši seka \emph{skyla},
	tai reiškia - egzistuoja homomorfizmas $\eta: C_n \to D_n$,
	toks kad $\xi \circ \eta = \mathup id_{C_n}$.
\end{proposition}
\begin{proof}
	Pasižymime
	\begin{equation}
		C_n = \present{a}{a^n=e}
		\,,\quad
		C_2 = \present{b}{b^2=e}
		\,.
	\end{equation}
	Turime atvaizdavimus
	\begin{equation}
		\begin{aligned}
			\zeta:	&& C_n &\to D_n \,,& a &\mapsto r \,,\\
			\xi:	&& D_n &\to C_2 \,,& r &\mapsto e \,, &s &\mapsto b\,,\\
			\eta:	&& C_2 &\to D_n \,,& b &\mapsto s \,.              
		\end{aligned}
	\end{equation}
	Atvizdavimams apibėžti užtenka pateikti,
	į ką siunčiami generatoriai.
	Kiti elementai grupėse yra žodžiai išreikšti generatorias
	ir yra atvaizduojami \textquote{paraidžiui}.
	Šie atvaizdavimai yra homomorfizmai,
	jei ir tik jei apibrėžiantis sąryšiai yra atvaizduojami į tapatybes
	atvaizdavimo kodomene.
	Parodome kad atvaizdavimai $\zeta$, $\xi$ ir $\eta$ yra grupės homomorfizmai
	\begin{equation}
		\begin{gathered}
			\zeta(a^n) = \zeta(a)^n = r^n = e = \zeta(e)
			\,,\\
			\xi(r^n) = \xi(r)^n = e^n = e = \xi(e)
			\,,\\
			\xi(s^2) = \xi(s)^2 = b^2 = e = \xi(e)
			\,,\\
			\xi(srs^{-1}) = \xi(s)\xi(r)\xi(s)^{-1} = b e b^{-1} = e = \xi(r^{-1})
			\,,\\
			\eta(b^2) = \eta(b)^2 = s^2 = e = \eta(e)
			\,.
		\end{gathered}
	\end{equation}
	test test

	Matome kad atvaizdavimo $\zeta$ vaizdas sutampa su $\xi$ branduoliu,
	tai reiškia $C_n$ yra normalus $D_n$ pogrupis.
	Taip pat matome, kad $\xi \circ \eta = \mathup id_{C_n}$.
	
	Nuskaitome homomorfizmą $\varphi: C_2 \to \mathup Aut(C_n)$,
	su kuriuo $C_2$ veikia į $C_n$
	\begin{equation}
		\begin{aligned}
			e &\mapsto && (& r &\mapsto ere^{-1} = r      &&)&\,,\\
			s &\mapsto && (& r &\mapsto srs^{-1} = r^{-1} &&)&\,.
		\end{aligned}
	\end{equation}

\end{proof}

\begin{definition}
	\textquote{Ciklinės grupės su vertikalia atspindžio polkštuma} yra 
	\begin{equation}
		C_{nv} := \present{r,h}{r^n=v^2=e,\, vrv^{-1}=r^{-1}}
		\,,\quad
		n \in \set N
		\,,
	\end{equation}
	čia generatorius $r$ yra tapatinamas su $2\pi/n$ pasukimu aplink ašį $z$,
	generatorius $v$ yra tapatinamas atspindžiu per vertikalią plokštumą,
	kuriai priklauso ašis $z$.
\end{definition}

\begin{corollary}
	Iš $C_{nv}$ apibėžimo iškart seka, kad $C_{nv} \cong D_n \cong C_n \ltimes C_2$. 
\end{corollary}

\begin{definition}
	\textquote{Ciklinės grupės su horizontalia atspindžio polkštuma} yra 
	\begin{equation}
		C_{nh} := \present{r,h}{r^n=h^2=e,\, hr=rh}
		\,,\quad
		n \in \set N
		\,,
	\end{equation}
	čia generatorius $r$ yra tapatinamas su $2\pi/n$ pasukimu aplink ašį $z$,
	generatorius $h$ yra tapatinamas atspindžiu per ašiai $z$ horizontalią plokštumą.
\end{definition}

\begin{proposition}
	\label{prop:cnh_product}
	Grupė $C_{nh}$ yra izomorfiška tiesioginiai grupių sandaugai $C_n \times C_2$
\end{proposition}
\begin{proof}
	Pasižymime $C_n = \present{a}{a^n=e}$, $C_2 = \present{b}{b^2=e}$
	Sąryšis $hr=rh$ leidžia kiek vieną $C_{nh}$ elementą užrašyti,
	forma $r^i h^j$, $i\in \qty{1,2,\dots, n}$, $j\in \qty{1,2}$.
	Tegul atvaizdavimas $f: C_{nh} \to C_n \times C_2$ yra 
	$r^i h^j \mapsto (a^i,h^j)$.
	Egzistuoja $f^{-1}: (a^i,h^j) \mapsto r^i h^j$.
	Apibėžiantis sąryšis $hr=rh$ atvaizduojamas į tapatybę grupėje $C_n \times C_2$
	\begin{equation}
		(a,e)(e,b) = (a,b) = (e,b)(a,e)
		\,.
	\end{equation}
	Tad atvaizdavimas $f$ yra izomorfizmas ir $C_{nh} \cong C_n \times C_2$
\end{proof}


\begin{definition}
	\textquote{Diiedrinės grupės su horizontalia atspindžio polkštuma} yra 
	\begin{equation}
		D_{nh} := \present{r,h}{
			r^n=s^2=h^2=e,\,
			hr=rh,\,
			hs=sh,\,
			srs^{-1}=r^{-1}
		}
		\,,\quad
		n \in \set N
		\,,
	\end{equation}
	čia generatorius $r$ yra tapatinamas su $2\pi/n$ pasukimu aplink ašį $z$,
	generatorius $h$ yra tapatinamas atspindžiu per ašiai $z$ horizontalią plokštumą,
	generatorius $s$ yra tapatinamas sukimu kampu $\pi$ aplink ašį statmeną ašiai $z$.
\end{definition}

\begin{proposition}
	Grupė $D_{nh}$ yra izomorfiška grupių tiesioginei sandaugai $D_n \times C_2$.
\end{proposition}
\begin{proof}
	Sąryšiai $hr=rh$ ir $hs=sh$ leidžia kiek vieną $D_{nh}$ elementą užrašyti,
	kaip $wh^j$, čia $j \in \qty{1,2}$ ir $w$ yra žodis užrašytas raidėmis $r$, $s$.
	Žodžiai $w$ turi tenkinti sąryšius $r^n=s^2=e$, $srs^{-1}=r^{-1}$,
	tad jie sudaro $D_{nh}$ pogrupį izomorfišką grupei $D_n$.
	Turint išskaidymą $wh^j$ įrodymas tesiamas taip pat,
	kaip ir teigynyje \ref{prop:cnh_product}.
\end{proof}


\begin{definition}
	\textquote{Atspindžių-posūkių grupės} (ang. \emph{the group of roto-reflections}) yra
	\begin{equation}
		S_{2n} := \present{f}{f^{2n}=e}
		\,,\quad
		n \in \set N
		\,,
	\end{equation}
	čia generatorius $f$ yra tapatinamas su netrikruoju pasukimu vadinamu atspindžio-posūkiu.
	Atspindys-posūkis yra sudarytas iš dviejų dalių:
	pirma pasukimo kampu $\pi/n$ aplink $z$ ašį,
	anta atspindžio per horizontalią plokštumą statmena $z$ ašiai.
	Ekinvalenčiai $S_{2n}$ galima apibėžti,
	kaip $C_{2nh}$ pogrupį generuojamą elemento $rh$.
\end{definition}

\begin{corollary}
	Grupė $S_{2n} \cong C_{2n}$.
\end{corollary}

\begin{definition}
	\textquote{Diiedrinės atspindžių-posūkių grupės} yra
	\begin{equation}
		D_{nd} := \present{f,v}{f^{2n}=e,\, vfv^{-1}=f^{-1}}
		\,,\quad
		n \in \set N
		\,,
	\end{equation}
	čia generatorius $f$ vėl yra tapatinamas su atspindžiu-posūkiu,
	generatorius $v$ yra tapatinamas su atspindžiu per vertikalią polkštūmą.
	Ekinvalenčiai $D_{nd}$ galima apibėžti,
	kaip $D_{nh}$ pogrupį generuojamą elemento $rsh$.
\end{definition}

\begin{corollary}
	Grupė $D_{nd} \cong C_{2n} \ltimes C_2$.
\end{corollary}

Grupės $C_n$ ir $D_n$ yra pogrupiai $\SO3$,
todėl jos vadinamos \emph{chiralinėmis} grupėmis.
Likusios ašinės grupės $C_{nv}$, $C_{nh}$, $D_{nh}$, $S_{2n}$ ir $D_nd$ nėra $\SO3$
pogrupiai, todėls yra vadinamos \emph{achiralinėmis} grupėmis.




\subsection{Poliedrinės grupės}

Poliedrinės grupės (ang. \emph{polyhedral groups}) yra $\Orth3$ pogrupiai,
kurie veikimu į $R^3$ fiksuoja poliedro (briaunainio) taškus.
Egzistuoja septynios poliedrinės grupės.

\begin{definition}
	\textquote{Chiralinė tetraedro simetrija} yra grupė 
	\begin{equation}
		T = \present{s,t}{s^2=t^3=(st)^3=e}
		\,,
	\end{equation}
	čia generatorius $s$ yra tapatinamas su sukimu kampu $\pi$ aplink ašį einančia per tetraedro briaunos vidurio tašką,
	generatorius $t$ yra tapatinamas su sukimu kampu $2\pi/3$ aplink ašį einančia per tetraedro sienos centrą,
	elementas $st$ yra tapatinamas su sukimu kampu $2\pi/3$ aplink ašį einančia per tetraedro viršūnę.
\end{definition}

\begin{definition}
	\textquote{Chiralinė oktaedro simetrija} yra grupė 
	\begin{equation}
		O = \present{s,t}{s^2=t^3=(st)^4=e}
		\,,
	\end{equation}
	čia generatorius $s$ yra tapatinamas su sukimu kampu $\pi$ aplink ašį einančia per oktaedro briaunos vidurio tašką,
	generatorius $t$ yra tapatinamas su sukimu kampu $2\pi/3$ aplink ašį einančia per oktaedro sienos centrą,
	elementas $st$ yra tapatinamas su sukimu kampu $2\pi/4$ aplink ašį einančia per oktaedro viršūnę.
\end{definition}

\begin{definition}
	\textquote{Chiralinė ikosaedro simetrija} yra grupė 
	\begin{equation}
		I = \present{s,t}{s^2=t^3=(st)^5=e}
		\,,
	\end{equation}
	čia generatorius $s$ yra tapatinamas su sukimu kampu $\pi$ aplink ašį einančia per ikosaedro briaunos vidurio tašką,
	generatorius $t$ yra tapatinamas su sukimu kampu $2\pi/3$ aplink ašį einančia per ikosaedro sienos centrą,
	elementas $st$ yra tapatinamas su sukimu kampu $2\pi/5$ aplink ašį einančia per ikosaedro viršūnę.
\end{definition}

\begin{definition}
	\textquote{Pilnoji tetraedro simetrija} yra grupė 
	\begin{equation}
		T_d = \present{s,t,v}{s^2=t^3=(st)^3=e,\, vsv^{-1}=s^1, vtv^{-1}=t^{-1}}
		\,
	\end{equation}
	čia elementai $s$, $t$ ir $st$ tapatinami taip pat kaip grupėje $T$,
	generatorius $v$ yra tapatinamas su atspindžiu per plokštumą,
	kuriai priklauso tetraedro briauna ir centras.
\end{definition}

\begin{definition}
	\textquote{Pilnoji oktaedro simetrija} yra grupė 
	\begin{equation}
		O_h = \present{s,t}{
			s^2=t^3=(st)^4=i^2=e,\,
			is=si ,\,
			it=ti }
		\,,
	\end{equation}
	čia elementai $s$, $t$ ir $st$ tapatinami taip pat kaip grupėje $O$,
	generatorius $i$ yra tapatinamas su inversija.
\end{definition}

\begin{definition}
	\textquote{Pilnoji ikosaedro simetrija} yra grupė 
	\begin{equation}
		I_h = \present{s,t}{
			s^2=t^3=(st)^5=i^2=e,\,
			is=si ,\,
			it=ti }
		\,,
	\end{equation}
	čia elementai $s$, $t$ ir $st$ tapatinami taip pat kaip grupėje $I$,
	generatorius $i$ yra tapatinamas su inversija.
\end{definition}

\begin{definition}
	\textquote{Piritoedro simetrija} yra grupė
	\begin{equation}
		T_h = \present{s,t}{
			s^2=t^3=(st)^3=i^2=e,\,
			is=si ,\,
			it=ti }
	\end{equation}
	čia elementai $s$, $t$ ir $st$ tapatinami taip pat kaip grupėje $T$,
	generatorius $i$ yra tapatinamas su inversija.
\end{definition} 

Grupės $T$, $O$ ir $I$ yra \emph{chiralinės grupės},
nes jos yra grupės $\SO3$ pogrupiai.
Grupės $T_d$, $O_h$, $I_h$ ir $T_h$ yra \emph{achiralinės grupės},
nes joms priklauso netikrieji pasukimai.



\clearpage{}
